% Formula sheet, Topics 1-2B — SHORT version (2 pages, two columns, 9pt). Condensed from formula-sheet-topics-1-2b.tex.
% Compile: pdflatex twice; check pdffonts (Type 1, uni = yes) and that the PDF has 2 pages.
\documentclass[10pt,a4paper]{extarticle}
\usepackage[utf8]{inputenc}
\usepackage[T1]{fontenc}
\usepackage{lmodern}   % Type 1 vector fonts (replaces PK bitmaps)
\usepackage{cmap}      % ToUnicode CMap -> copyable text and accents
\usepackage[english]{babel}
\usepackage{amsmath,amssymb}
\usepackage[a4paper,margin=0.9cm,top=0.8cm,bottom=0.9cm]{geometry}
\usepackage{multicol}
\usepackage{enumitem}
\usepackage{xcolor}
\usepackage{microtype}
\setlist{nosep,leftmargin=0.9em,label=\textbullet}
\pagestyle{empty}
\setlength{\parindent}{0pt}
\setlength{\parskip}{1.2pt}
\setlength{\columnsep}{9pt}
\setlength{\columnseprule}{0.2pt}
\setlength{\abovedisplayskip}{1.5pt}\setlength{\belowdisplayskip}{1.5pt}
\setlength{\abovedisplayshortskip}{1pt}\setlength{\belowdisplayshortskip}{1pt}
\newcommand{\T}[1]{\par\smallskip{\color{blue!55!black}\bfseries\large #1}\par\smallskip}
\newcommand{\B}[1]{\par\medskip{\color{blue!55!black}\bfseries #1}\par}
\newcommand{\E}{\mathbb{E}}
\newcommand{\dd}{\displaystyle}

\begin{document}
\begin{center}
{\bfseries\large Formula sheet --- Topics 1, 2A, 2B}\quad {\small Finan\c{c}as II (MPE Insper)}\\[-1pt]
{\scriptsize $V=D+E$; $\tau$ corporate tax; $r_D,r_E,r_A$ expected returns; $r_U$ unlevered cost of capital; $g$ growth; $N$ shares.}
\end{center}
\vspace{-6pt}
\begin{multicols}{2}

% ===================================================================== TOPIC 1
\T{Topic 1 --- Accounting}
\B{Balance sheet}
$A=L+E$ (book and market); $\Delta A=\Delta L+\Delta E$.\\
Debt $\min(X,F)$ + equity $\max(X-F,0)=X$; $FCF_A=FCF_L+FCF_E$.\\
$NWC=CA-CL$;\ Dep $=(K_0-S)/n$;\ $BV=$ cost $-$ acc.\ dep.\\
Goodwill $=$ price $-$ fair value of net identifiable assets.
\begin{itemize}
\item Liabilities senior/specific, equity junior/residual. Intangibles mostly off the balance sheet; Dep is not cash.
\end{itemize}

\B{Book vs.\ market, leverage, liquidity}
$E^{mkt}=PN$;\ $M/B=E^{mkt}/E^{book}$;\ net debt $=D-$Cash;\\
$EV=E^{mkt}+D-\text{Cash}$;\ current $=CA/CL$;\ quick $=(CA-\text{Inv})/CL$.
\begin{itemize}
\item Market forward-, book backward-looking; book $E$ can be $<0$. High $M/B$ = growth. Covenants on book ratios.
\end{itemize}

\B{Income statement}
Cash collected $=S-\Delta AR$;\ $GP=S-\text{COGS}$;\ $EBIT=GP-\text{Opex}$;\\
$NI=(EBIT-\text{Int})(1-\tau)$;\ $E_{end}=E_{beg}+NI-\text{Div}$.\\
$OI=(1-v)S-F$;\ $DOL=\%\Delta OI/\%\Delta S=1+F/OI$.
\begin{itemize}
\item Accrual and matching principles; Dep spreads capex. COGS $\approx$ variable, opex $\approx$ fixed.
\end{itemize}

\B{Ratios}
$EPS=NI/N$; margins $GP/S$, $OI/S$, $NI/S$; AR days $=AR/(S/365)$; AP days $=AP/(\text{COGS}/365)$; coverage $=EBIT/\text{Int}$; $P/E$; economic earnings $=NI-r_EE$.
\[ROE=\tfrac{NI}{E}=\tfrac{NI}{S}\cdot\tfrac{S}{A}\cdot\tfrac{A}{E}\ \text{(DuPont)},\quad ROA=\tfrac{NI}{A}\]
\[ROIC=\tfrac{EBIT(1-\tau)}{D+E},\quad ROE=ROIC+\tfrac DE[ROIC-r_D(1-\tau)]\]
$\operatorname{sd}(ROE)=(1+D/E)\operatorname{sd}(ROIC)$.
\begin{itemize}
\item Leverage raises ROE iff $ROIC>r_D(1-\tau)$, always raises risk. ROE ignores timing, risk, uses book values.
\end{itemize}

\B{Cash flows and FCF}
$CFO=NI+\text{Dep}+\text{non-cash}-\Delta AR+\Delta AP-\Delta\text{Inv}$;\ $\Delta\text{Cash}=CFO+CFI+CFF$.
\[\boxed{FCF=EBIT(1-\tau)+\text{Dep}-\text{Capex}-\Delta NWC}\]
\[\boxed{FCFE=FCF-\text{Int}(1-\tau)+\Delta D}\]
$CF_D=\text{Int}-\Delta D$;\ $FCFE+CF_D=FCF+\tau\,\text{Int}$.
\begin{itemize}
\item FCF $\to$ EV (WACC, or $r_U$ in APV); FCFE $\to$ equity ($r_E$). $\Delta NWC$: operating items only. Add back non-cash charges incl.\ deferred taxes.
\end{itemize}

% ===================================================================== TOPIC 2A
\T{Topic 2A --- Valuation, NPV, MM}
\B{Time value of money}
Perpetuity $C/r$;\ growing $C_1/(r-g)$;\ annuity $\frac Cr[1-(1+r)^{-T}]$;\\
growing annuity $\frac{C_1}{r-g}[1-(\frac{1+g}{1+r})^T]$;\ $EAR=(1+\frac{APR}{k})^k-1$;\\
$1+r_{nom}=(1+r_{real})(1+\pi)$.

\B{Present value and bonds}
\[P_A=\sum_t\tfrac{\E[\tilde d_t]}{(1+r_t)^t},\quad P=\sum_{t\le T}\tfrac{c}{(1+r_t)^t}+\tfrac{F}{(1+r_T)^T}\]
Zero: $r_T=(F/P)^{1/T}-1$;\ YTM $y$: $P=\sum CF_t/(1+y)^t$;\\
$(1+r_2)^2=(1+r_1)(1+f_{1,2})$;\ $HPR=(P_1-P_0+c)/P_0$;\\
$D_{mac}=\frac1P\sum t\,\frac{CF_t}{(1+y)^t}$,\ $\frac{\Delta P}{P}\approx-\frac{D_{mac}}{1+y}\Delta y$.
\begin{itemize}
\item Price $\downarrow$ when $y\uparrow$; premium iff $c>y$; pulls to par. Long, low-coupon = more sensitive. Callable $=$ bond $-$ call. Spread $=y_{corp}-y_{gov}$.
\end{itemize}

\B{CAPM, betas, comparables}
$r=r_f+\beta(\E[r_m]-r_f)$;\ $\beta=\text{Cov}(r,r_m)/\text{Var}(r_m)$;
\[r_A=\tfrac EV r_E+\tfrac DV r_D,\quad \beta_A=\tfrac EV\beta_E+\tfrac DV\beta_D\]
$r_E=r_A+\frac DE(r_A-r_D)$;\ $\beta_E=\beta_A+\frac DE(\beta_A-\beta_D)$.
\begin{itemize}
\item Use the project's risk. Unlever each comparable, average $\beta_A$, relever at target $D/E$.
\end{itemize}

\B{Which cash flows; NPV; IRR}
After-tax salvage $=SV-\tau(SV-BV)$;\ $TV_T=FCF_{T+1}/(r-g)$;\\
$NPV=PV-c$;\ $IRR$: $0=-c+\sum \E[d_t]/(1+IRR)^t$.
\begin{itemize}
\item Incremental: opportunity costs and side effects in; sunk costs and unchanged overhead out; financing not in FCF; NWC invested then recovered; nominal with nominal.
\item NPV $>0$ $\Leftrightarrow$ shareholder gain (any financing, frictionless). Fisher separation.
\item IRR pitfalls: borrowing-type flows, multiple IRRs, scale/timing (use NPV or crossover rate).
\end{itemize}

\B{Debt vs.\ equity; state prices}
$C_D=\min(A,F)=F-\max(F-A,0)$ (bond $-$ put);\ $C_E=\max(A-F,0)$ (call).\\
Price $=\sum_s p_sx_s$;\ $\sum_sp_s=\frac1{1+r_f}$;\ $\pi_s=p_s(1+r_f)$;\ price $=\frac{\sum\pi_sx_s}{1+r_f}$;\\
put--call parity $C-P=S-K/(1+r_f)$.
\begin{itemize}
\item Benchmark: competitive complete markets, same borrowing terms, no taxes, no bankruptcy costs, cash flows independent of financing. No arbitrage $\Leftrightarrow p_s>0$; find $p_s$ by replication.
\end{itemize}

\B{MM I, MM II, cost of debt}
$V_L=V_U$;\ $V(K)=C(K)+\frac{K}{1+r_f}-P(K)=S$.
\[r_E=r_A+\tfrac DE(r_A-r_D),\qquad r_{WACC}=r_A=r_U\ \text{(no taxes)}\]
$D=\frac{F}{1+y}=\frac{F(1-p)+Cp}{1+r_D}\Rightarrow y>r_D$ if $p>0$.
\begin{itemize}
\item Proof by arbitrage with homemade leverage (or undoing it). Conservation of value: repackaging risk adds nothing.
\item $r_E$ linear in $D/E$ while debt is safe, flatter once $r_D$ rises. Yield $=$ promised, not expected, return.
\end{itemize}

% ===================================================================== TOPIC 2B
\T{Topic 2B --- Taxes, distress, trade-off}
\B{Tax shield and MM with taxes}
$ITS=\tau r_DD=\tau\cdot\text{Int}$;\ payout $(1-\tau)X+\tau r_DD$;
\[\boxed{V_L=V_U+PV(ITS)}=D+E\]
$\frac DVr_D+\frac EVr_E=\frac{V_U}{V_L}r_U+\frac{V_{ITS}}{V_L}r_{ITS}$.\\
Recap: $P=\frac{V_U+PV(ITS)}{N_0}$ at announcement; shares bought $=D/P$.

\B{By debt policy}
\textit{Constant debt} ($r_{ITS}=r_D$): $V_{ITS}=\tau D$;\ $r_E=r_U+\frac DE(1-\tau)(r_U-r_D)$;\ $\beta_E=\beta_U[1+(1-\tau)\frac DE]$;\ $r_{WACC}=r_U(1-\tau\frac DV)$ ($g=0$).\\
\textit{Target $D/V$, continuous} ($r_{ITS}=r_U$): $V_{ITS}=\frac{\tau r_DD}{r_U-g}$;\ $r_E=r_U+\frac DE(r_U-r_D)$;\ $r_{WACC}=r_U-\tau r_D\frac DV$.\\
\textit{Annual}: multiply $V_{ITS}$ by $\frac{1+r_U}{1+r_D}$;\ $r_{WACC}=r_U-\tau r_D\frac DV\frac{1+r_U}{1+r_D}$.

\B{WACC and APV}
\[\boxed{r_{WACC}=\tfrac EVr_E+\tfrac DVr_D(1-\tau)}\]
\[r_{WACC}-g=(r_U-g)\big(1-\tfrac{V_{ITS}}{V_L}\big)\]
$APV=V_U+PV(ITS)-PV(\text{distress})+PV(\text{side effects})$.
\begin{itemize}
\item WACC: unlever comparables, relever at target $D/E$; assumes constant $D/E$. APV: known debt schedule. Same answer under constant $D/E$.
\item Risky debt: shield on coupon paid, $\bar r_D=\E[C_{t+1}]/D_t$; after-tax cost $r_D-\tau\bar r_D$.
\end{itemize}

\B{Puzzle, personal taxes, distress}
$V_L/V_U=1+\tau D/V_U$;\ Miller $G_L=\big[1-\frac{(1-\tau_c)(1-\tau_s)}{1-\tau_b}\big]D$.
\begin{itemize}
\item Default = transfer of control, not an MM violation. Direct costs (legal, fire sales) vs.\ indirect (customers, suppliers, workers; asset substitution, overhang, cashing out). Borne by shareholders ex ante.
\end{itemize}

\B{Static trade-off}
$V_U=\frac{(1-\tau)\E[X]}{1+r_f}$;\ safe debt ($K<\underline X$): $V=V_U+\frac{\tau K}{1+r_f}$;
\[V(K)=V_U+\tfrac{1}{1+r_f}\Big[\tau\Big(\textstyle\int_{\underline X}^{K}xh+K(1-H)\Big)-CH(K)\Big]\]
FOC $\tau[1-H(K^*)]=Ch(K^*)$;\ $dK^*/d\tau>0$, $dK^*/dC<0$;\\
uniform: $K^*=\bar X-C/\tau$ (if $C<\tau(\bar X-\underline X)$).

\B{Bank capital; payout}
\begin{itemize}
\item More equity: ROE falls but so does risk and required $r_E$ (MM II); lost shield is a transfer, not a social cost; failure externality $\Rightarrow$ social optimum has more equity.
\end{itemize}
Sources $=$ uses: $X_t+\text{issues}_t=I_t+\text{Div}_t$;\ $V_t=\E\sum_{i\ge1}\delta^i(X_{t+i}-I_{t+i})$.
\begin{itemize}
\item Payout irrelevant given investment; dividend $\equiv$ repurchase in a perfect market.
\end{itemize}

\T{Exam traps}
\begin{itemize}
\item Negative book equity $\ne$ insolvency; profit $\ne$ cash; Dep is not a source of cash (only its tax saving is).
\item EV $=E^{mkt}+D-$Cash, not market cap $+$ total liabilities. Say which $D/E$ (book or market).
\item FCF taxes EBIT as if unlevered: no interest, no tax shield in FCF (shield goes in WACC or APV).
\item $\Delta NWC$: operating items only (not excess cash, not short-term debt).
\item Higher ROE $\ne$ better: check $ROIC$ vs.\ $r_D(1-\tau)$ and the added risk.
\item Use asset betas of comparables; never discount a project at the firm's $r_E$.
\item ``Debt is cheaper'' is wrong without taxes: $r_E$ rises to offset (MM II), WACC unchanged.
\item Equity gets riskier with leverage even when debt is riskless (more claims on the same risk).
\item The yield (or coupon) is not $r_D$ when default is possible.
\item State prices are not probabilities ($\sum p_s=1/(1+r_f)$; $\pi_s$ are risk-neutral, not real).
\item IRR does not rank projects (scale, timing, sign changes): use NPV.
\item $V_{ITS}=\tau D$ only for permanent, constant, safe debt; with a target $D/E$ use $r_U$.
\item $r_{WACC}=r_U(1-\tau D/V)$ only for constant debt and $g=0$.
\item Bankruptcy costs, not default itself, break MM; they are paid by shareholders ex ante.
\end{itemize}

\end{multicols}
\end{document}
